%==============================================================================
%  config/packages.tex  --  Paquetes y configuracion base
%  Tesis doctoral - Universidad Nacional de San Luis (UNSL)
%------------------------------------------------------------------------------
%  Preambulo modernizado y deduplicado (junio 2026).
%  Cambios clave respecto del template original:
%    - Se ELIMINO el conflicto color/xcolor: ahora se carga SOLO xcolor.
%    - Se quitaron cargas DUPLICADAS (graphicx, amsmath, setspace, color...).
%    - Se removieron paquetes innecesarios para esta tesis (epsfig, CJKutf8,
%      fontenc LGR/griego).
%    - Se agregaron microtype y booktabs (tipografia y tablas modernas).
%    - hyperref se carga casi al final, como corresponde, y cleveref despues.
%    - Los paquetes de algoritmos se cargan de forma condicional.
%==============================================================================

% ---- Codificacion e idioma --------------------------------------------------
\usepackage[utf8]{inputenc}      % archivos fuente en UTF-8 (acentos, n)
\usepackage[T1]{fontenc}         % salida con acentos correctos / copiar-pegar
\usepackage[spanish,es-tabla,es-noindentfirst,es-nodecimaldot]{babel}
% es-tabla: los flotantes dicen "Tabla"; es-noindentfirst: sin sangria inicial.

% ---- Fuentes ----------------------------------------------------------------
% Tipografia: Computer Modern (la que usa la tesis de referencia de la UNSL).
% A proposito NO se carga lmodern, para que la fuente coincida con la del
% director (con T1 + cm-super salen las mismas SFRM*; en TeX Live completo y
% Overleaf ya estan). Si alguna vez quisieras la version moderna, agrega:
%   \usepackage{lmodern}
\usepackage{courier}             % monoespaciada para codigo en linea
% microtype: protrusion (mejora margenes). expansion=false para no depender de
% fuentes escalables (cm-super) al usar Computer Modern -> compila en todos lados.
\usepackage[expansion=false]{microtype}

% ---- Matematica -------------------------------------------------------------
\usepackage{amsmath}
\usepackage{amssymb}
\usepackage{amsfonts}
\usepackage{amsthm}
\usepackage{amsxtra}
\usepackage{mathtools}           % extiende amsmath (\DeclarePairedDelimiter...)

% ---- Color: UN SOLO paquete (xcolor reemplaza por completo a color) ----------
\usepackage[dvipsnames,svgnames,table]{xcolor}
% 'table' habilita colores de fila/celda; 'dvipsnames'/'svgnames' dan nombres.
% NO cargar 'color' por separado: ese era el origen de los errores de choque.

% Colores institucionales UNSL (del escudo), para acentos del template:
\definecolor{unslazul}{RGB}{48,48,144}       % azul del escudo
\definecolor{unslceleste}{RGB}{120,192,240}  % celeste del escudo

% ---- Figuras y tablas -------------------------------------------------------
\usepackage{graphicx}
\graphicspath{{images/}}         % raiz para \includegraphics (mas subcarpetas)
\usepackage{rotating}            % figuras/tablas apaisadas
% rac.sty redefine \@makecaption con un formato propio que el paquete caption
% no reconoce (warning "Unknown document class") y que igualmente pisa con sus
% defaults. Restauramos la definicion estandar de LaTeX antes de cargarlo para
% que la reconozca; el resultado visual es el mismo (manda caption).
\makeatletter
\long\def\@makecaption#1#2{%
  \vskip\abovecaptionskip
  \sbox\@tempboxa{#1: #2}%
  \ifdim \wd\@tempboxa >\hsize
    #1: #2\par
  \else
    \global \@minipagefalse
    \hb@xt@\hsize{\hfil\box\@tempboxa\hfil}%
  \fi
  \vskip\belowcaptionskip}
\makeatother
\usepackage{subcaption}          % subfiguras (reemplaza al viejo 'subfigure')
\usepackage[font={small,bf}]{caption}
\usepackage{booktabs}            % \toprule \midrule \bottomrule
\usepackage{multirow}
\usepackage{makecell}
\usepackage{adjustbox}           % escalar/ajustar tablas y cajas anchas

% ---- Espaciado, notas, varios -----------------------------------------------
\usepackage{setspace}            % \onehalfspacing, \doublespacing
\usepackage[bottom]{footmisc}    % notas al pie siempre al fondo de la pagina
\usepackage{verbatim}
\usepackage{soul}                % \hl (resaltar), \st (tachar)
\usepackage{etoolbox}            % \patchcmd y utilidades de programacion

% ---- Pagina: A4, margenes simetricos y encabezado ---------------------------
% Reemplaza los margenes "Rackham" de rac.sty por A4 con margenes simetricos
% (igual izquierda/derecha) y mas aire arriba/abajo. Agrega un encabezado con
% el nombre del capitulo y el numero de pagina, con una linea debajo.
\usepackage[a4paper,left=3cm,right=3cm,top=3cm,bottom=3cm,
            headheight=15pt,headsep=0.7cm]{geometry}
\usepackage{fancyhdr}
\pagestyle{fancy}
\fancyhf{}                                  % limpia encabezado y pie
\fancyhead[L]{\nouppercase{\leftmark}}      % capitulo a la izquierda
\fancyhead[R]{\thepage}                     % numero de pagina a la derecha
\renewcommand{\headrulewidth}{0.4pt}        % linea bajo el encabezado
\renewcommand{\footrulewidth}{0pt}
% Aperturas de capitulo: sin encabezado, numero de pagina abajo-centro
\fancypagestyle{plain}{%
  \fancyhf{}%
  \fancyfoot[C]{\thepage}%
  \renewcommand{\headrulewidth}{0pt}%
}

% ---- Algoritmos (carga condicional) -----------------------------------------
% algorithm + algorithmicx/algpseudocode estan en TeX Live completo y Overleaf.
% Se cargan solo si estan disponibles para no romper entornos minimos.
\IfFileExists{algpseudocode.sty}{%
  \usepackage{algorithm}%
  \usepackage{algpseudocode}%
}{%
  \PackageWarning{packages}{algpseudocode no encontrado: se omiten los entornos de algoritmos.}%
}

% ---- Codigo fuente: minted --------------------------------------------------
% minted requiere compilar con  -shell-escape  y tener Pygments instalado
% (pip install Pygments). En Overleaf funciona sin configuracion extra.
% Alternativa sin dependencias: comentar esta linea y usar 'listings'.
% La opcion outputdir ya no es necesaria (ni valida) con minted v3+ / TeX Live 2024+:
% el directorio de salida se detecta automaticamente.
\usepackage{minted}

% ---- Bibliografia -----------------------------------------------------------
\usepackage{natbib}              % citas \citep / \citet (estilo en thesis.tex)
% Por defecto se usa plainnat. El template trae agu04.bst en styles/ como
% alternativa (estilo de la AGU); para usarlo: \bibliographystyle{agu04}.
% El encabezado de la bibliografia lo pone \startbibliography (estilo capitulo);
% anulamos el encabezado propio de natbib para que no quede duplicado:
\renewcommand{\bibsection}{}

% ---- Abreviaturas / acronimos -----------------------------------------------
\usepackage[printonlyused]{acronym}   % lista de abreviaturas (lee abbr.tex)

% ---- Hiperenlaces (cargar casi al final) ------------------------------------
\usepackage[spanish,unicode]{hyperref}
\hypersetup{
  colorlinks=true,
  linkcolor=unslazul,         % enlaces internos (capitulos, secciones)
  citecolor=unslazul,         % citas bibliograficas
  urlcolor=unslazul,          % URLs
  linktoc=all,                % en el indice, enlaza numero y texto
}
\renewcommand{\UrlFont}{\color{unslazul}\rmfamily}

% ---- Referencias cruzadas inteligentes (despues de hyperref) ----------------
\usepackage[capitalise,noabbrev]{cleveref}   % \cref{...} -> "figura 3.1", etc.

% rac.sty usa un tipo de entrada 'chap' en el indice; le indicamos a hyperref
% su nivel de marcador para evitar el warning "bookmark level for unknown chap".
\makeatletter
\providecommand*{\toclevel@chap}{0}
\makeatother

% ---- Nota sobre estilos heredados del template ------------------------------
% styles/aas_macros.sty  : macros de revistas astronomicas (ADS). NO se carga
%   porque no es necesario en esta tesis; queda por compatibilidad. Si lo
%   necesitaras: \usepackage{aas_macros}
% styles/agu04.bst        : estilo bibliografico de la AGU (geofisica), opcional.
% styles/rac.sty          : formato de tesis exigido (UNSL/Rackham); SE USA.

